% =====================================================================
% TEXTO ORIGINAL DO CAPÍTULO DE FUNDAMENTAÇÃO TEÓRICA — 19 páginas
% Salvo em 2026-04-30 como referência para comparação com V3-SEGURA
% =====================================================================

\chapter{FUNDAMENTAÇÃO TEÓRICA}

{\color{blue} A fundamentação teórica representa a base epistemológica e metodológica desta investigação, pois explicita os referenciais que orientam tanto a análise do problema em estudo quanto a elaboração do artefato didático proposto. No âmbito da Educação Matemática, construir um referencial teórico vai além da simples apresentação de autores; envolve compreender de que modo distintas concepções de conhecimento, ensino e aprendizagem repercutem em decisões didáticas concretas (Artigue, 2014; Brousseau, 1997).

O referencial teórico exerce função estruturante ao contribuir para a delimitação do problema, a definição dos objetivos e a seleção dos procedimentos metodológicos, posicionando o objeto investigado em um quadro mais amplo. No contexto desta pesquisa, a fundamentação teórica orienta a elaboração de uma sequência didática digital destinada ao ensino de Probabilidade no Ensino Médio, articulando princípios da Engenharia Didática, da Teoria das Situações Didáticas, dos ciclos iterativos de concepção em ambientes digitais e do construcionismo.

A opção pelo método Design Science Research (Dresch; Lacerda; Antunes Júnior, 2015) pressupõe que o artefato desenvolvido esteja ancorado em bases teóricas consistentes, uma vez que a intervenção pedagógica é concebida como resposta fundamentada a uma classe de problemas educacionais. Assim, a fundamentação teórica assume papel orientador e normativo, oferecendo critérios para a organização da sequência didática, a seleção dos conteúdos de Probabilidade e a integração de recursos digitais.

Dessa forma, esta seção estrutura-se em três eixos principais: (i) a discussão das sequências didáticas como forma racional de organização do ensino; (ii) os fundamentos do ensino de Probabilidade no Ensino Médio; e (iii) o papel dos Objetos Virtuais de Aprendizagem na mediação do conhecimento matemático. O primeiro eixo é apresentado a seguir.}

% [TEXTO INTEGRAL DAS SEÇÕES PRESERVADO NA MENSAGEM ORIGINAL DO USUÁRIO,
%  DATA 2026-04-30. SUBSEÇÕES MAPEADAS:
%   - Sequências Didáticas (introdução)
%   - Abordagens Metodológicas para a Construção de Sequências Didáticas
%   - Engenharia Didática
%   - Teoria das Situações Didáticas
%   - Particularidades das abordagens com OVAs
%   - Como elaborar uma Sequência Didática no Ensino de Matemática
%   - Progressão conceitual
%   - Problematização
%   - Coerência
%   - Papel da mediação docente na aprendizagem
%   - Sequências Didáticas para o Ensino de Probabilidade: BNCC e recursos digitais
% ]
%
% NOTA: o texto integral está disponível na transcrição da sessão de
% 2026-04-30. Este arquivo é stub de referência; o conteúdo integral
% foi enviado pelo orientando como input do prompt e está preservado no
% histórico da sessão. A V3-SEGURA opera sobre essa versão como base.
